\documentclass[11pt,oneside]{article}
\usepackage[a4paper,margin=26mm,headheight=15pt]{geometry}
\usepackage[T1]{fontenc}
\usepackage[utf8]{inputenc}
\IfFileExists{libertinus.sty}{\usepackage{libertinus}}{\usepackage{lmodern}}
\usepackage{microtype}
\usepackage{amsmath,amssymb,amsthm,mathtools,mathrsfs,bm}
\usepackage{booktabs,longtable,array,tabularx,multirow}
\usepackage{graphicx}
\usepackage{xcolor}
\usepackage{enumitem}
\usepackage{fancyhdr}
\usepackage{titlesec}
\usepackage{listings}
\usepackage{xurl}
\usepackage{hyperref}
\usepackage[nameinlink,noabbrev]{cleveref}

\definecolor{linkblue}{RGB}{25,75,135}
\definecolor{shadegray}{RGB}{246,247,249}
\definecolor{rulegray}{RGB}{195,201,208}
\definecolor{deepgreen}{RGB}{30,105,75}
\hypersetup{
  colorlinks=true,
  pdfauthor={OpenAI, prepared for Vladimir Reshetnikov},
  linkcolor=linkblue,
  citecolor=deepgreen,
  urlcolor=linkblue,
  pdftitle={Fabius Monomial Antiderivatives and Inverse-Quantile Integrals},
  pdfsubject={Exact primitive ladders, complex-exponent series, Rvachev transforms, and inverse Fabius integrals},
  pdfkeywords={Fabius function, inverse Fabius function, Rvachev up-function, antiderivative, Mellin transform, Thue-Morse, q-binomial, Bell polynomial, Bernoulli number, Lambert W}
}

\pagestyle{fancy}
\fancyhf{}
\fancyhead[L]{\small Fabius monomial antiderivatives}
\fancyhead[R]{\small\nouppercase{\leftmark}}
\fancyfoot[C]{\thepage}
\renewcommand{\headrulewidth}{0.4pt}
\titleformat{\section}{\Large\bfseries}{\thesection}{0.7em}{}
\titleformat{\subsection}{\large\bfseries}{\thesubsection}{0.7em}{}
\titleformat{\subsubsection}{\normalsize\bfseries}{\thesubsubsection}{0.7em}{}
\setcounter{secnumdepth}{3}
\setcounter{tocdepth}{2}
\setlist[itemize]{leftmargin=2em,itemsep=0.25em,topsep=0.4em}
\setlist[enumerate]{leftmargin=2.3em,itemsep=0.25em,topsep=0.4em}
\setlength{\emergencystretch}{3em}
\allowdisplaybreaks[2]
\numberwithin{equation}{section}

\newtheorem{theorem}{Theorem}[section]
\newtheorem{proposition}[theorem]{Proposition}
\newtheorem{lemma}[theorem]{Lemma}
\newtheorem{corollary}[theorem]{Corollary}
\newtheorem{identity}[theorem]{Identity}
\newtheorem{conjecture}[theorem]{Conjecture}
\theoremstyle{definition}
\newtheorem{definition}[theorem]{Definition}
\newtheorem{example}[theorem]{Example}
\newtheorem{algorithm}[theorem]{Algorithm}
\newtheorem{obligation}[theorem]{Formalization obligation}
\theoremstyle{remark}
\newtheorem{remark}[theorem]{Remark}
\newtheorem{warning}[theorem]{Warning}

\newcommand{\R}{\mathbb R}
\newcommand{\C}{\mathbb C}
\newcommand{\Q}{\mathbb Q}
\newcommand{\N}{\mathbb N}
\newcommand{\Z}{\mathbb Z}
\newcommand{\E}{\mathbb E}
\newcommand{\Prob}{\mathbb P}
\newcommand{\Up}{\operatorname{up}}
\newcommand{\sinc}{\operatorname{sinc}}
\newcommand{\Mellin}{\mathcal M}
\newcommand{\Vol}{\mathcal J}
\newcommand{\InvF}{J}
\newcommand{\e}{\mathrm e}
\newcommand{\ii}{\mathrm i}
\newcommand{\dd}{\,\mathrm d}
\newcommand{\bigO}{\mathcal O}
\newcommand{\smallo}{o}
\newcommand{\supp}{\operatorname{supp}}
\newcommand{\wt}{\operatorname{wt}_2}
\newcommand{\fall}[2]{#1^{\underline{#2}}}
\newcommand{\rise}[2]{#1^{\overline{#2}}}
\newcommand{\abs}[1]{\left|#1\right|}
\newcommand{\norm}[1]{\left\lVert#1\right\rVert}
\newcommand{\pospart}[1]{\left(#1\right)_{+}}
\newcommand{\qbinom}[3]{\genfrac{[}{]}{0pt}{}{#1}{#2}_{#3}}
\newcommand{\repo}{\nolinkurl{Analysis/FabiusFunction/docs}}
\newcommand{\decl}[1]{\nolinkurl{#1}}
\newcommand{\familyname}[1]{\nolinkurl{#1}}
\newcommand{\newresult}{\textbf{Derived here relative to the audited repository snapshot.}}
\newcommand{\corpusresult}{\textbf{Imported from the repository corpus.}}
\newcommand{\conjectural}{\textbf{Conjectural.}}

\newenvironment{boxedremark}{%
  \par\smallskip\noindent\begin{center}\begin{minipage}{0.94\textwidth}%
  \color{black}\hrule\vspace{0.6em}\small
}{%
  \vspace{0.6em}\hrule\end{minipage}\end{center}\smallskip
}

\lstdefinestyle{pseudo}{
  basicstyle=\ttfamily\small,
  backgroundcolor=\color{shadegray},
  frame=single,
  rulecolor=\color{rulegray},
  columns=fullflexible,
  keepspaces=true,
  showstringspaces=false,
  breaklines=true,
  xleftmargin=0.7em,
  xrightmargin=0.7em,
  aboveskip=0.8em,
  belowskip=0.8em
}

\title{\bfseries Fabius Monomial Antiderivatives and Inverse-Quantile Integrals\\[0.45em]
\large Exact primitive ladders, complex-exponent series, Rvachev transforms,\\
Mellin interpolation, and nonlinear inverse reductions}
\author{Frontier research report prepared from the recursive TeX corpus\\
\small Vladimir Reshetnikov's \texttt{ProveIt} repository}
\date{27 August 2026\\
\small Snapshot-relative novelty audit; no claim of unconditional worldwide priority}

\begin{document}
\maketitle

\begin{abstract}
Let $F$ be the bounded Fabius distribution function, let $\Up(x)=F(1-|x|)$ be
Rvachev's compactly supported even up-function, let $\mathcal F$ be the signed
global extension satisfying $\mathcal F'(x)=2\mathcal F(2x)$, and let
$\InvF=F^{-1}$ on $[0,1]$.  A recursive audit of the LaTeX documents under
\repo{} finds extensive treatments of exact dyadic arithmetic, moments,
$q$-binomial and Thue--Morse formulae, repeated-integration limits, Fourier sinc
products, endpoint asymptotics, and inverse asymptotics, but no systematic
calculus of indefinite monomial-weighted primitives.

The central observation is the exact zero-based primitive ladder
\[
 A_n(x)=2^{\binom n2}F(x/2^n),\qquad A_n'=A_{n-1},\qquad A_0=F,
\]
valid for $0\le x\le1$.  It immediately gives, for every $p\in\N_0$,
\[
 \int x^pF(x)\dd x=
 \sum_{k=0}^{p}(-1)^k\frac{p!}{(p-k)!}
 2^{\binom{k+1}{2}}x^{p-k}F(x/2^{k+1})+C.
\]
For arbitrary $\alpha\in\C$ the terminating sum becomes an absolutely and
locally uniformly convergent series.  Endpoint flatness is strong enough that
this series is exact for every complex exponent, including all negative
powers, rather than merely asymptotic.  Consequently
$\Mellin_F(s)=\int_0^1x^{s-1}F(x)\dd x$ is entire and has an explicit dyadic
Newton-type expansion.  Differentiation in $s$ produces logarithmic weights,
with coefficients expressible through signed Stirling numbers or complete Bell
polynomials.

Reflection gives primitives of $F(1-x)=1-F(x)$, and the same formulas provide
piecewise primitives and all polynomial moments of $\Up$.  For the inverse
function a change of variables yields the exact nonlinear reduction
\[
 \int_0^y t^\beta\InvF(t)\dd t=
 \frac{y^{\beta+1}\InvF(y)}{\beta+1}
 -\frac1{\beta+1}\int_0^{\InvF(y)}F(x)^{\beta+1}\dd x,
\]
with a separate logarithmic identity at $\beta=-1$.  Positive integer powers of
$F$ in the remainder admit an explicit finite Thue--Morse spline limit.
The repository's Lambert--$W$ inverse asymptotic then gives sharp endpoint laws
for these inverse primitives, including a phase transition at $\beta=-1$.
Exact rational checks and independent quasi-Monte Carlo experiments accompany
the proofs.  Several conjectural extensions and a concrete Lean formalization
roadmap conclude the report.
\end{abstract}

\tableofcontents
\newpage

\section{Executive summary and nonduplication boundary}
\label{sec:executive}

\subsection{Corpus audit}
The audit covered the living expositions, the consolidated frontier volume, the
paper transcriptions, the glossary and Lean walkthrough, the currently visible
frontier-draft families, and every frozen historical family listed by the
repository archive manifest.  The tree is actively edited, so the audit was
performed at the document-family level: principal mathematical sources were
read closely, draft reports were searched thematically and lexically, and
auxiliary generated tables or frozen duplicates were checked through their
parent documents and manifests.  The principal interfaces are the current
synthesis \cite{ProveItPrimary}, the consolidated frontier volume
\cite{ProveItFrontiers}, the inverse-frontier report \cite{ProveItInverse}, and
the archive ledger \cite{ProveItArchive}.

The closest existing material is the repeated-integration dossier in the
frontier volume.  It studies how normalized repeated integrations of splines
converge to $\Up$ and gives exact convergence rates.  Its object is a sequence
of approximating functions.  The present report instead fixes the limiting
Fabius function and identifies its entire Volterra primitive tower.  Searches
for \emph{antiderivative}, \emph{indefinite integral}, and the corresponding
monomial patterns returned no treatment of the formulas below in the current
primary synthesis, consolidated frontier volume, or the new inverse, Newton,
and spectral frontier reports.

\begin{boxedremark}
All novelty labels in this report mean ``not found in the recursively audited
repository snapshot of 27 August 2026.''  They are not claims that no equivalent
identity has ever appeared elsewhere.  Elementary consequences are included
because they organize a useful calculus; the frontier claims concern the
combined exact ladder, the all-complex-exponent convergence theorem, its Mellin
and Bell consequences, and the inverse nonlinear reduction with its
Thue--Morse realization.
\end{boxedremark}

\subsection{Contribution map}
\begin{table}[ht]
\centering
\small
\begin{tabularx}{\textwidth}{@{}p{0.25\textwidth}X p{0.17\textwidth}@{}}
\toprule
Topic & Result & Status \\
\midrule
Primitive ladder & $n$-fold zero-based primitive $2^{\binom n2}F(x/2^n)$ and its Volterra-kernel form. & Proved here \\
Polynomial weights & Finite closed form for $\int x^pF(x)\dd x$ and for arbitrary polynomial weights. & Proved here \\
Complex powers & Absolutely convergent exact series for every $\alpha\in\C$, with an exact remainder and compact-uniform convergence. & Proved here \\
Mellin calculus & Entire Mellin transform, complete monotonicity on the real axis, dyadic interpolation, and logarithmic derivatives. & Proved here \\
Arithmetic bridge & Complete moments as finite combinations of inverse-dyadic values, hence direct inputs for Bernoulli, Thue--Morse, and $q$-binomial formulae. & Proved here \\
Rvachev transforms & Piecewise antiderivatives, global normalized primitives, and fractional absolute moments. & Proved here \\
Inverse Fabius & Exact weighted reduction to $\int F^{\beta+1}$, critical logarithmic formula, convergence threshold, and order-statistic interpretation. & Proved here \\
Nonlinear remainder & Explicit finite Thue--Morse spline formula for $\int F^r$, $r\in\N$. & Proved here \\
Endpoint inverse & Weighted inverse-primitive equivalents, including the critical $\beta=-1$ law. & Proved from corpus asymptotic \\
Fractional primitive index & A continuous-index $q$-Volterra semigroup interpolating the integer ladder. & Conjectural program \\
\bottomrule
\end{tabularx}
\caption{Main contributions and status.}
\label{tab:contribution-map}
\end{table}

\section{Established platform and notation}
\label{sec:platform}

\subsection{The bounded Fabius function}
We use the repository's normalization.  The bounded Fabius function
$F:\R\to[0,1]$ satisfies
\begin{equation}
 F(x)=0\quad(x\le0),\qquad F(x)=1\quad(x\ge1),             \label{eq:F-exterior}
\end{equation}
\begin{equation}
 F'(x)=2F(2x)\quad(0\le x\le\tfrac12),\qquad
 F(1-x)=1-F(x)\quad(x\in\R).                              \label{eq:F-basic}
\end{equation}
It is $C^\infty$, strictly increasing on $[0,1]$, and flat at both endpoints.
Let
\begin{equation}
 X=\sum_{j\ge1}2^{-j}U_j,
 \qquad U_j\stackrel{\mathrm{iid}}{\sim}\operatorname{Unif}[0,1]. \label{eq:X-random}
\end{equation}
Then $F$ is the distribution function of $X$.  In particular, for every
$z\in\C$, all complex moments $\E X^z$ exist after choosing the positive-real
branch, because the law has infinite-order flatness at zero.

Write
\begin{equation}
 d_n=\E X^n
     =n\int_0^1t^{n-1}\Up(t)\dd t
     =n!\,2^{\binom n2}F(2^{-n}),\qquad n\in\N,             \label{eq:d-moments}
\end{equation}
with $d_0=1$.  The repository proves the triangular recurrence
\begin{equation}
 (n+1)(2^n-1)d_n=\sum_{k=0}^{n-1}\binom{n+1}{k}d_k,         \label{eq:d-recurrence}
\end{equation}
and the Bernoulli recurrence
\begin{equation}
 d_n=\frac{n2^n}{4(2^n-1)}d_{n-1}
 -\frac1{2^n-1}\sum_{k=1}^{\lfloor n/2\rfloor}
 \binom n{2k}2^{n-2k}B_{2k}d_{n-2k}.                       \label{eq:d-bernoulli}
\end{equation}
Thus any result expressed through the $d_n$ is automatically connected to the
corpus's Bernoulli and exact dyadic layers.

\subsection{Rvachev's up-function and the signed extension}
The even compactly supported bump is
\begin{equation}
 \Up(x)=F(1-|x|),\qquad \supp\Up=[-1,1].                   \label{eq:up-def}
\end{equation}
Equivalently, $\Up$ is the density of $Y=2X-1$.  The signed global extension is
\begin{equation}
 \mathcal F(x)=\sum_{m\ge0}(-1)^{\wt(m)}\Up(x-2m-1),       \label{eq:global-def}
\end{equation}
a locally finite sum satisfying
\begin{equation}
 \mathcal F'(x)=2\mathcal F(2x)\qquad(x\in\R).             \label{eq:global-diffeq}
\end{equation}
It agrees with $F$ on $(-\infty,1]$ but is not the bounded distribution
function to the right of $1$.  This distinction is essential when extending
primitive formulas globally.

\subsection{Inverse function}
Let
\begin{equation}
 \InvF=F^{-1}:[0,1]\longrightarrow[0,1].                  \label{eq:J-def}
\end{equation}
Then $\InvF(1-y)=1-\InvF(y)$ and $\InvF$ is smooth on $(0,1)$, with infinite
endpoint slope.  The corpus establishes the leading lower-endpoint equivalent
\begin{equation}
 \InvF(y)\sim
 \frac{\sqrt{T}}{\e\sqrt{\log2}}
 \exp\!\left(-\sqrt{2(\log2)T}\right),
 \qquad T=\log\frac1y\longrightarrow\infty.               \label{eq:J-leading}
\end{equation}
The sharper repository theory contains a periodic Gamma--zeta correction and
Lambert--$W$ reversion; \eqref{eq:J-leading} is the interface needed for the
first endpoint laws below.

\section{The exact primitive ladder}
\label{sec:ladder}

\subsection{A self-similar Volterra chain}
For $n\in\N_0$ define
\begin{equation}
 A_0(x)=F(x),\qquad
 A_n(x)=2^{\binom n2}F(x/2^n)\quad(n\ge1).                 \label{eq:An-def}
\end{equation}

\begin{theorem}[Primitive ladder]
\label{thm:primitive-ladder}
For $0\le x\le1$ and every $n\ge1$,
\begin{equation}
 A_n'(x)=A_{n-1}(x),\qquad A_n(0)=0.                       \label{eq:ladder-derivative}
\end{equation}
Consequently $A_n$ is the $n$-fold zero-based primitive of $F$:
\begin{equation}
 \boxed{
 2^{\binom n2}F(x/2^n)
 =\frac1{(n-1)!}\int_0^x(x-t)^{n-1}F(t)\dd t.}            \label{eq:volterra-ladder}
\end{equation}
For the signed extension $\mathcal F$, the same statements hold for every
$x\in\R$.
\end{theorem}

\begin{proof}
For $n\ge1$, $0\le x/2^n\le1/2$, so \eqref{eq:F-basic} gives
\[
 A_n'(x)=2^{\binom n2-n}F'(x/2^n)
 =2^{\binom n2-n+1}F(x/2^{n-1})
 =2^{\binom{n-1}{2}}F(x/2^{n-1})=A_{n-1}(x).
\]
All $A_n(0)$ vanish.  Iterating the fundamental theorem of calculus gives the
Volterra kernel.  For $\mathcal F$, equation \eqref{eq:global-diffeq} has no
domain restriction, so the calculation is global.
\end{proof}

\begin{remark}[Why the triangular exponent is inevitable]
Let $(\Vol f)(x)=\int_0^xf(t)\dd t$ and $(Sf)(x)=f(x/2)$.  Then
\begin{equation}
 \Vol S=2S\Vol,\qquad \Vol F=SF.                           \label{eq:q-weyl}
\end{equation}
Therefore
\begin{equation}
 \Vol^nF=2^{0+1+\cdots+(n-1)}S^nF
 =2^{\binom n2}S^nF.                                      \label{eq:q-weyl-ladder}
\end{equation}
The coefficients are thus a $q$-Weyl phenomenon with $q=1/2$, not an accidental
by-product of repeated differentiation.  This is the operator-level bridge to
the half-base $q$-binomial calculus in the corpus.
\end{remark}

\subsection{Repeated integration by parts}
\begin{theorem}[Finite weighted identity]
\label{thm:finite-weighted}
Let $g\in C^N([0,x])$, where $0\le x\le1$.  Then
\begin{align}
 \int_0^xg(t)F(t)\dd t
 &=\sum_{k=0}^{N-1}(-1)^kg^{(k)}(x)A_{k+1}(x)              \notag\\
 &\quad+(-1)^N\int_0^xg^{(N)}(t)A_N(t)\dd t.              \label{eq:finite-weighted}
\end{align}
The identity is global with $F,A_n$ replaced by $\mathcal F,
2^{\binom n2}\mathcal F(\,\boldsymbol{\cdot}/2^n)$.
\end{theorem}

\begin{proof}
Use $A_{k+1}'=A_k$ and integrate by parts successively.  Every lower boundary
term vanishes because $A_k(0)=0$.
\end{proof}

\section{Closed forms for nonnegative integer powers}
\label{sec:polynomial}

\subsection{The monomial formula}
\begin{theorem}[Polynomial-weight antiderivative]
\label{thm:monomial}
For every $p\in\N_0$ and $0\le x\le1$,
\begin{equation}
 \boxed{
 \int x^pF(x)\dd x=
 \sum_{k=0}^{p}(-1)^k\frac{p!}{(p-k)!}
 2^{\binom{k+1}{2}}x^{p-k}F(x/2^{k+1})+C.}                \label{eq:monomial-closed}
\end{equation}
For the signed global extension $\mathcal F$, the same formula is valid for all
$x\in\R$.
\end{theorem}

\begin{proof}
Apply \cref{thm:finite-weighted} with $g(t)=t^p$ and $N=p+1$.  Since
$g^{(p+1)}=0$, the remainder vanishes.  Adding an arbitrary constant converts
the zero-based primitive to an indefinite antiderivative.
\end{proof}

The first cases are
\begin{align}
 \int F(x)\dd x
 &=F(x/2)+C,                                                \label{eq:p0}\\
 \int xF(x)\dd x
 &=xF(x/2)-2F(x/4)+C,                                      \label{eq:p1}\\
 \int x^2F(x)\dd x
 &=x^2F(x/2)-4xF(x/4)+16F(x/8)+C,                          \label{eq:p2}\\
 \int x^3F(x)\dd x
 &=x^3F(x/2)-6x^2F(x/4)+48xF(x/8)-384F(x/16)+C.            \label{eq:p3}
\end{align}
The cancellations are exact: differentiating the $k$th term produces one
ordinary monomial derivative and one dilation derivative; the latter cancels
the former contribution from the preceding scale.

\begin{corollary}[Arbitrary polynomial weight]
\label{cor:polynomial}
For a polynomial $P$ of degree $m$,
\begin{equation}
 \boxed{
 \int P(x)F(x)\dd x=
 \sum_{k=0}^{m}(-1)^k2^{\binom{k+1}{2}}
 P^{(k)}(x)F(x/2^{k+1})+C.}                                \label{eq:polynomial-closed}
\end{equation}
\end{corollary}

\subsection{Continuation outside the unit interval}
Formula \eqref{eq:monomial-closed} is a bounded-$F$ identity on $[0,1]$.
For $x>1$, define the normalized primitive
\begin{equation}
 I_p(x)=\int_0^xt^pF(t)\dd t.                              \label{eq:Ip-def-int}
\end{equation}
Then
\begin{equation}
 I_p(x)=I_p(1)+\frac{x^{p+1}-1}{p+1},\qquad x\ge1,         \label{eq:Ip-outside}
\end{equation}
while $I_p(x)=0$ for $x\le0$.  Confusing this piecewise continuation with the
global signed-extension formula is a normalization error: $F$ is constant for
$x\ge1$, whereas $\mathcal F$ continues in alternating compact blocks.

\subsection{Complete moments and exact arithmetic}
Putting $x=1$ in \eqref{eq:monomial-closed} gives a finite inverse-dyadic
formula.

\begin{corollary}[Moment bridge]
\label{cor:moment-bridge}
For $p\in\N_0$,
\begin{align}
 \int_0^1x^pF(x)\dd x
 &=\sum_{k=0}^{p}(-1)^k\frac{p!}{(p-k)!}
 2^{\binom{k+1}{2}}F(2^{-k-1})                             \label{eq:moment-dyadic}\\
 &=\frac1{p+1}\sum_{j=1}^{p+1}(-1)^{j-1}
   \binom{p+1}{j}d_j                                      \label{eq:moment-d}\\
 &=\boxed{\frac{1-d_{p+1}}{p+1}}.                         \label{eq:moment-prob}
\end{align}
In particular, every integer-weight complete integral is rational.
\end{corollary}

\begin{proof}
The first equality is evaluation of the primitive.  Use
$2^{\binom j2}F(2^{-j})=d_j/j!$ for the second.  Finally, integration by parts
against the probability measure $\dd F$ gives
\[
 d_{p+1}=\int_0^1x^{p+1}\dd F(x)
 =1-(p+1)\int_0^1x^pF(x)\dd x.
\]
\end{proof}

Thus the antiderivative calculus is a direct consumer of every exact formula
for $F(2^{-n})$ in the repository.  For example, with $V_n=F(2^{-n})$,
\begin{equation}
 V_0=1,\qquad
 V_n=\frac{2^{-\binom n2}}{2^n-1}
 \sum_{k=0}^{n-1}\frac{2^{\binom k2}}{(n-k+1)!}V_k,       \label{eq:V-recurrence}
\end{equation}
so \eqref{eq:moment-dyadic} is immediately executable in exact rational
arithmetic.  Alternatively one may insert the corpus's finite Thue--Morse
formula or the translated $q$-binomial formula below.  For
$m,n\in\N$ and arbitrary $Q\in\C$, its right-hand side is independent of
$Q$ (and equals the bounded value $F(m/2^n)$ whenever $m\le2^n$):
\begin{align}
 \mathcal F(m/2^n)
 &=\frac{1}{2^{n^2}(1/2;1/2)_n}
 \sum_{k=0}^{n}
 \frac{\qbinom nk{1/2}}{4^{\binom k2}(n+k)!}              \notag\\
 &\quad\times\sum_{r=0}^{m2^k-1}(-1)^{\wt(r)}
 (r-m2^k+Q)^{n+k},                                        \label{eq:q-dyadic-import}
\end{align}
or the Bernoulli recurrence \eqref{eq:d-bernoulli}.  No new summation theorem
is required: the primitive formula provides the missing finite linear
functional that transports those arithmetic representations into integration.

The first values are
\begin{equation}
\begin{array}{c|cccccc}
 p&0&1&2&3&4&5\\ \hline
 \displaystyle\int_0^1x^pF(x)\dd x
 &\frac12&\frac{13}{36}&\frac5{18}&\frac{1207}{5400}
 &\frac{251}{1350}&\frac{22661}{142884}.
\end{array}                                                \label{eq:first-moments-table}
\end{equation}

\subsection{A scaled reflection identity}
Because $F(1-x)=1-F(x)$, the same integrand has two different primitive
representations.  Equating them gives a family of functional identities.

\begin{proposition}[Integrated reflection]
\label{prop:integrated-reflection}
Let $I_j(x)=\int_0^xt^jF(t)\dd t$.  For $p\in\N_0$ and $0\le x\le1$,
\begin{equation}
 \boxed{
 \sum_{j=0}^{p}(-1)^j\binom pj I_j(1-x)
 =I_p(x)-\frac{x^{p+1}}{p+1}+\frac{d_{p+1}}{p+1}.}         \label{eq:integrated-reflection}
\end{equation}
After inserting \eqref{eq:monomial-closed}, this becomes a finite identity
between values of $F$ at the two dyadic scale towers
$x/2^r$ and $(1-x)/2^r$.
\end{proposition}

\begin{proof}
Substitution $u=1-x$ gives
\[
 \int x^pF(1-x)\dd x
 =-\sum_{j=0}^{p}(-1)^j\binom pjI_j(1-x)+C.
\]
The same integrand is $x^p(1-F(x))$, with primitive
$x^{p+1}/(p+1)-I_p(x)$.  Determine the constant at $x=1$ using
\eqref{eq:moment-prob}.
\end{proof}

\section{Fractional, negative, and complex powers}
\label{sec:complex-powers}

\subsection{Flatness estimate}
The all-complex-exponent result rests on more than formal repeated integration.

\begin{lemma}[Uniform Volterra flatness]
\label{lem:volterra-flatness}
For every integer $s\ge0$ there is a constant $C_s$ such that, for all
$n\ge1$ and $0\le x\le1$,
\begin{equation}
 0\le A_n(x)\le C_s\frac{s!}{(s+n)!}x^{s+n}
 =C_s\frac{x^{s+n}}{(s+1)_n}.                              \label{eq:An-flat-bound}
\end{equation}
\end{lemma}

\begin{proof}
Endpoint flatness gives $F(t)\le C_st^s$ on $[0,1]$.  Insert this estimate into
\eqref{eq:volterra-ladder} and evaluate the beta integral:
\[
 A_n(x)\le\frac{C_s}{(n-1)!}\int_0^x(x-t)^{n-1}t^s\dd t
 =C_s\frac{s!}{(s+n)!}x^{s+n}.
\]
\end{proof}

\subsection{The exact complex-exponent series}
For $x>0$ use the principal real logarithm and write $x^\alpha=
\exp(\alpha\log x)$.  Define
\begin{equation}
 I_\alpha(x)=\int_0^xt^\alpha F(t)\dd t.                  \label{eq:Ialpha-def}
\end{equation}
Unlike the corresponding integral with $1-F$, this converges for every
$\alpha\in\C$, because $F$ is flat at zero.

\begin{theorem}[All-complex-exponent antiderivative]
\label{thm:complex-series}
For every $\alpha\in\C$ and $0<x\le1$,
\begin{equation}
 \boxed{
 I_\alpha(x)=\sum_{k=0}^{\infty}(-1)^k\fall{\alpha}{k}
 2^{\binom{k+1}{2}}x^{\alpha-k}F(x/2^{k+1}).}              \label{eq:complex-series}
\end{equation}
The series is absolutely convergent and locally uniformly convergent in
$(\alpha,x)\in\C\times(0,1]$.  It therefore defines an entire function of
$\alpha$.  For $\alpha=p\in\N_0$ it terminates at $k=p$ and reduces to
\eqref{eq:monomial-closed}.

For every $N\ge1$ the exact remainder after $N$ terms is
\begin{equation}
 R_N(\alpha,x)=(-1)^N\fall{\alpha}{N}
 \int_0^xt^{\alpha-N}A_N(t)\dd t.                          \label{eq:complex-remainder}
\end{equation}
\end{theorem}

\begin{proof}
Apply \cref{thm:finite-weighted} with $g(t)=t^\alpha$.  This gives the first
$N$ terms and \eqref{eq:complex-remainder}.  Choose $s$ so that
$\Re\alpha+s+1>0$.  By \cref{lem:volterra-flatness},
\begin{equation}
 \abs{R_N(\alpha,x)}
 \le \frac{C_s\abs{\fall{\alpha}{N}}}{(s+1)_N}
 \frac{x^{\Re\alpha+s+1}}{\Re\alpha+s+1}.                \label{eq:R-bound}
\end{equation}
For nonintegral $\alpha$,
$\fall{\alpha}{N}=(-1)^N\Gamma(N-\alpha)/\Gamma(-\alpha)$, so the quotient
in \eqref{eq:R-bound} is $\bigO(N^{-\Re\alpha-s-1})$.  For nonnegative integer
$\alpha$ it vanishes identically for large $N$.  Uniform gamma-ratio estimates
on compact $\alpha$-sets, with the polynomial definition of the falling
factorial used across its removable singularities, give compact-uniform
convergence.  The same bound applied termwise proves absolute convergence.
Entire dependence under the integral also follows directly from dominated
convergence, using flatness of arbitrarily high order.
\end{proof}

\begin{corollary}[One-step scale recurrence]
\label{cor:Ialpha-scale}
For every $\alpha\in\C$ and $0<x\le1$,
\begin{equation}
 \boxed{
 I_\alpha(x)=x^\alpha F(x/2)-\alpha2^\alpha
 I_{\alpha-1}(x/2).}                                      \label{eq:Ialpha-scale}
\end{equation}
Here $2^\alpha=\exp(\alpha\log2)$.  Iterating
\eqref{eq:Ialpha-scale} gives the finite expansion and exact remainder in
\cref{thm:complex-series}; letting the depth tend to infinity gives
\eqref{eq:complex-series}.
\end{corollary}

\begin{proof}
Integrate once by parts with the primitive $F(t/2)$, and substitute $t=2u$ in
the remaining integral.  Convergence at the lower endpoint holds for every
$\alpha$ by flatness.
\end{proof}

\begin{remark}[A convergent integration-by-parts expansion]
Ordinary repeated integration by parts often produces a divergent asymptotic
series.  Here the rapidly increasing dilation factor
$2^{\binom{k+1}{2}}$ is exactly defeated by the much smaller value
$F(x/2^{k+1})$.  The Volterra beta-factor $(s+1)_{k+1}^{-1}$ exposes the
mechanism.  The result is a genuine convergent expansion for every complex
exponent.
\end{remark}

For $x>1$ the bounded-function continuation is
\begin{equation}
 I_\alpha(x)=I_\alpha(1)+
 \begin{cases}
 \dfrac{x^{\alpha+1}-1}{\alpha+1},&\alpha\ne-1,\\[6pt]
 \log x,&\alpha=-1,
 \end{cases}                                               \label{eq:Ialpha-outside}
\end{equation}
where the first quotient has its usual removable value at $\alpha=-1$.  For
$\mathcal F$ the series \eqref{eq:complex-series}, with $F$ replaced by
$\mathcal F$, is valid for every $x>0$ because the dilation equation is global.

\subsection{Negative integer powers}
When $\alpha=-m$, $m\in\N$, the two signs cancel because
$\fall{-m}{k}=(-1)^k\rise{m}{k}$.

\begin{corollary}[Positive-term negative-power series]
\label{cor:negative-series}
For $m\in\N$ and $0<x\le1$,
\begin{equation}
 \boxed{
 \int_0^x\frac{F(t)}{t^m}\dd t
 =\sum_{k=0}^{\infty}\rise{m}{k}
 2^{\binom{k+1}{2}}x^{-m-k}F(x/2^{k+1}).}                 \label{eq:negative-series}
\end{equation}
All terms are nonnegative.  In particular,
\begin{equation}
 \int_0^x\frac{F(t)}t\dd t
 =\sum_{k=0}^{\infty}k!\,2^{\binom{k+1}{2}}
 x^{-k-1}F(x/2^{k+1}).                                    \label{eq:negative-one-series}
\end{equation}
\end{corollary}

This positivity gives certified lower bounds from partial sums.  Combining
\eqref{eq:R-bound} with an explicit, formally verified flatness constant
would turn it into a fully explicit rational enclosure at dyadic $x$.

\subsection{Logarithmic powers, Stirling numbers, and Bell polynomials}
Since $I_\alpha(x)$ is entire,
\begin{equation}
 \frac{\partial^r}{\partial\alpha^r}I_\alpha(x)
 =\int_0^xt^\alpha(\log t)^rF(t)\dd t.                    \label{eq:log-derivative}
\end{equation}
Termwise differentiation in \eqref{eq:complex-series} gives
\begin{align}
 \int_0^xt^\alpha(\log t)^rF(t)\dd t
 &=\sum_{k\ge0}(-1)^k2^{\binom{k+1}{2}}F(x/2^{k+1})x^{\alpha-k} \notag\\
 &\quad\times\sum_{\ell=0}^{r}\binom r\ell
 (\log x)^{r-\ell}\frac{\dd^\ell}{\dd\alpha^\ell}
 \fall{\alpha}{k}.                                       \label{eq:log-series}
\end{align}
The derivative may be written without singularities through signed Stirling
numbers $s(k,j)$:
\begin{equation}
 \frac{\dd^\ell}{\dd\alpha^\ell}\fall{\alpha}{k}
 =\sum_{j=\ell}^{k}s(k,j)\fall{j}{\ell}\alpha^{j-\ell}.  \label{eq:stirling-derivative}
\end{equation}
Away from the removable roots $0,1,\ldots,k-1$, it also has the Bell form
\begin{equation}
 \frac{\dd^\ell}{\dd\alpha^\ell}\fall{\alpha}{k}
 =\fall{\alpha}{k}\,
 \mathcal B_\ell(h_1,\ldots,h_\ell),                     \label{eq:bell-derivative}
\end{equation}
where $\mathcal B_\ell$ is the complete exponential Bell polynomial and
\begin{equation}
 h_j=\psi^{(j-1)}(\alpha+1)-\psi^{(j-1)}(\alpha-k+1).      \label{eq:bell-h}
\end{equation}
This supplies a direct antiderivative connection to the Bell-polynomial
machinery used elsewhere in the corpus for saddle and inverse expansions.

At $\alpha=0$, $\frac{\dd}{\dd\alpha}\fall{\alpha}{k}|_{\alpha=0}
=(-1)^{k-1}(k-1)!$ for $k\ge1$, so
\begin{equation}
 \boxed{
 \int_0^xF(t)\log t\dd t
 =F(x/2)\log x-
 \sum_{k=1}^{\infty}(k-1)!\,2^{\binom{k+1}{2}}
 x^{-k}F(x/2^{k+1}).}                                     \label{eq:log-explicit}
\end{equation}

\section{The Mellin transform as an entire dyadic interpolant}
\label{sec:mellin}

Define
\begin{equation}
 \Mellin_F(s)=\int_0^1x^{s-1}F(x)\dd x.                   \label{eq:Mellin-def}
\end{equation}

\begin{theorem}[Entire Mellin series]
\label{thm:mellin-entire}
The function $\Mellin_F$ is entire and
\begin{equation}
 \boxed{
 \Mellin_F(s)=\sum_{k=0}^{\infty}(-1)^k\fall{s-1}{k}
 2^{\binom{k+1}{2}}F(2^{-k-1}).}                          \label{eq:Mellin-series}
\end{equation}
For $\Re s>0$,
\begin{equation}
 \Mellin_F(s)=\frac{1-\E X^s}{s}.                        \label{eq:Mellin-prob}
\end{equation}
Since $s\mapsto\E X^s$ is entire, this quotient continues to every
$s\in\C$; its singularity at $s=0$ is removable, with
\begin{equation}
 \Mellin_F(0)=-\E\log X.                                 \label{eq:Mellin-zero}
\end{equation}
\end{theorem}

\begin{proof}
Use \cref{thm:complex-series} with $\alpha=s-1$.  For $\Re s>0$, integration
by parts against $\dd F$ gives
$\E X^s=1-s\Mellin_F(s)$.  Both sides are entire, so the identity
continues to all $s$; l'Hopital's rule gives \eqref{eq:Mellin-zero}.
\end{proof}

\begin{proposition}[Strict complete monotonicity and log-convexity]
\label{prop:mellin-cm}
For every real $s$ and $r\in\N_0$,
\begin{equation}
 (-1)^r\Mellin_F^{(r)}(s)
 =\int_0^1x^{s-1}(-\log x)^rF(x)\dd x>0.                  \label{eq:Mellin-cm}
\end{equation}
Consequently $\Mellin_F$ is strictly completely monotone on the entire real
axis and strictly log-convex there.
\end{proposition}

\begin{proof}
Flatness justifies differentiation for every real $s$.  Positivity is strict
because the integrand is positive on a set of positive measure.  Log-convexity
follows from Cauchy--Schwarz applied to the measure
$x^{s-1}F(x)\dd x$.
\end{proof}

The change of variables $x=\e^{-u}$ gives the bilateral-strength Laplace
representation
\begin{equation}
 \Mellin_F(s)=\int_0^\infty\e^{-su}F(\e^{-u})\dd u,       \label{eq:Mellin-Laplace}
\end{equation}
which converges for every real $s$ because $F(\e^{-u})$ decays faster than every
ordinary exponential.  Thus the inverse-dyadic values in
\eqref{eq:Mellin-series} interpolate an entire, positive, log-convex Mellin
object rather than merely its integer samples.

\section{\texorpdfstring{The transforms $1-F(x)$ and $F(1-x)$}{The transforms 1-F(x) and F(1-x)}}
\label{sec:transforms}

The bounded normalization satisfies $F(1-x)=1-F(x)$ on all of $\R$.
For $x>0$ and $\alpha\ne-1$, an indefinite primitive is
\begin{equation}
 \boxed{
 \int x^\alpha(1-F(x))\dd x
 =\frac{x^{\alpha+1}}{\alpha+1}-I_\alpha(x)+C.}            \label{eq:complement-primitive}
\end{equation}
At $\alpha=-1$ it is
\begin{equation}
 \int\frac{1-F(x)}x\dd x=\log x-I_{-1}(x)+C.             \label{eq:complement-minus-one}
\end{equation}
The normalized integral from zero exists only for $\Re\alpha>-1$, because
$1-F(0)=1$; this contrasts with $I_\alpha$, which exists for every complex
$\alpha$.

For $p\in\N_0$, substitution $u=1-x$ gives the second finite form
\begin{equation}
 \boxed{
 \int x^pF(1-x)\dd x
 =-\sum_{j=0}^{p}(-1)^j\binom pj I_j(1-x)+C.}              \label{eq:reflected-primitive}
\end{equation}
Equations \eqref{eq:complement-primitive} and
\eqref{eq:reflected-primitive} are equal up to the constant computed in
\cref{prop:integrated-reflection}.

For example,
\begin{align}
 \int F(1-x)\dd x&=x-F(x/2)+C,                             \label{eq:reflect-p0}\\
 \int xF(1-x)\dd x&=\frac{x^2}{2}-xF(x/2)+2F(x/4)+C.      \label{eq:reflect-p1}
\end{align}

\section{Rvachev up-function primitives}
\label{sec:up}

\subsection{Piecewise polynomial weights}
On $0<x<1$, $\Up(x)=1-F(x)$, so
\begin{equation}
 \int x^p\Up(x)\dd x=\frac{x^{p+1}}{p+1}-I_p(x)+C.        \label{eq:up-positive}
\end{equation}
On $-1<x<0$, $\Up(x)=F(1+x)$ and, for integer $p$,
\begin{equation}
 \int x^p\Up(x)\dd x
 =\sum_{j=0}^{p}(-1)^{p-j}\binom pj I_j(1+x)+C.           \label{eq:up-negative}
\end{equation}

\begin{theorem}[Global normalized polynomial primitive of $\Up$]
\label{thm:up-global-primitive}
Let
\begin{equation}
 H_p(x)=\int_{-\infty}^{x}t^p\Up(t)\dd t.                 \label{eq:Hp-def}
\end{equation}
Then
\begin{equation}
 H_p(x)=
 \begin{cases}
 0,&x\le-1,\\[2pt]
 \displaystyle\sum_{j=0}^{p}(-1)^{p-j}\binom pjI_j(1+x),
   &-1\le x\le0,\\[8pt]
 \displaystyle\frac{(-1)^pd_{p+1}}{p+1}
 +\frac{x^{p+1}}{p+1}-I_p(x),&0\le x\le1,\\[9pt]
 \displaystyle\frac{1+(-1)^p}{p+1}d_{p+1},&x\ge1.
 \end{cases}                                               \label{eq:Hp-piecewise}
\end{equation}
\end{theorem}

\begin{proof}
The first two branches follow from \eqref{eq:up-negative}.  Evenness gives
\[
 \int_{-1}^{0}t^p\Up(t)\dd t
 =(-1)^p\int_0^1t^p\Up(t)\dd t
 =\frac{(-1)^pd_{p+1}}{p+1}.
\]
Add the positive-side primitive \eqref{eq:up-positive}.  The final branch is
the complete moment.
\end{proof}

For $p=0$ this gives the particularly simple distribution function of the
up-density:
\begin{equation}
 H_0(x)=
 \begin{cases}
 0,&x\le-1,\\
 F((1+x)/2),&-1\le x\le0,\\
 \frac12+x-F(x/2),&0\le x\le1,\\
 1,&x\ge1.
 \end{cases}                                               \label{eq:up-cdf}
\end{equation}

\subsection{Fractional and negative powers}
For $0<r\le1$ define the positive-side complement primitive
\begin{equation}
 \Phi_\alpha(r)=
 \begin{cases}
 \dfrac{r^{\alpha+1}}{\alpha+1}-I_\alpha(r),&\alpha\ne-1,\\[7pt]
 \log r-I_{-1}(r),&\alpha=-1.
 \end{cases}                                               \label{eq:Phi-alpha}
\end{equation}
Then $\Phi_\alpha'(r)=r^\alpha\Up(r)$ on $(0,1)$.  Hence for every complex
$\alpha$ it is a local antiderivative on the positive half of the support.  If
on $-1<x<0$ one chooses the upper-half-plane branch
$x^\alpha=\e^{\ii\pi\alpha}(-x)^\alpha$, then
\begin{equation}
 \boxed{
 \int x^\alpha\Up(x)\dd x=
 \begin{cases}
 -\e^{\ii\pi\alpha}\Phi_\alpha(-x)+C,&-1<x<0,\\[3pt]
 \Phi_\alpha(x)+C,&0<x<1.
 \end{cases}}                                             \label{eq:up-complex-local}
\end{equation}
The lower-half-plane branch replaces $\e^{\ii\pi\alpha}$ by
$\e^{-\ii\pi\alpha}$.  Since $I_\alpha$ is given by
\eqref{eq:complex-series}, these are explicit convergent series on both open
half-intervals, including every negative and fractional exponent.

A single primitive across the origin exists exactly when $\Re\alpha>-1$ for
the branch-free weight $|x|^\alpha$.  Put
\begin{equation}
 T_\alpha=\Phi_\alpha(1)
 =\frac{\E X^{\alpha+1}}{\alpha+1}.                       \label{eq:T-alpha}
\end{equation}
Then the global normalized primitive is
\begin{equation}
 \boxed{
 \int_{-\infty}^{x}|t|^\alpha\Up(t)\dd t=
 \begin{cases}
 0,&x\le-1,\\
 T_\alpha-\Phi_\alpha(-x),&-1\le x\le0,\\
 T_\alpha+\Phi_\alpha(x),&0\le x\le1,\\
 2T_\alpha,&x\ge1.
 \end{cases}}                                             \label{eq:up-absolute-primitive}
\end{equation}
In particular,
\begin{equation}
 \boxed{
 \int_{-1}^{1}|x|^\alpha\Up(x)\dd x
 =\frac{2\E X^{\alpha+1}}{\alpha+1},\qquad \Re\alpha>-1.} \label{eq:up-fractional-moment}
\end{equation}
The right side supplies a meromorphic continuation to all $\alpha\in\C$, with
a simple pole of residue $2$ at $\alpha=-1$.  The pole records the genuine
nonintegrability of $|x|^{-1}\Up(x)$ at the central point, in contrast with
the endpoint-flat integral $I_{-1}$.

For $m\in\N_0$,
\begin{equation}
 \int_{-1}^{1}x^{2m}\Up(x)\dd x
 =\frac{2d_{2m+1}}{2m+1},\qquad
 \int_{-1}^{1}x^{2m+1}\Up(x)\dd x=0.                     \label{eq:up-integer-moments}
\end{equation}
Thus the antiderivative formula recovers the moment bridge between the Fabius
distribution and the infinite sinc product
$\widehat\Up(\xi)=\prod_{j\ge1}\sinc(2^{-j}\xi)$: the Taylor coefficients of
the product are exactly the complete even moments in
\eqref{eq:up-integer-moments}.

\section{The signed global extension}
\label{sec:global}

Since \eqref{eq:global-diffeq} holds on all of $\R$, the cleanest version of the
calculus belongs to $\mathcal F$.

\begin{theorem}[Global polynomial and complex-power formulas]
\label{thm:global-formulas}
For $p\in\N_0$ and all $x\in\R$,
\begin{equation}
 \int x^p\mathcal F(x)\dd x
 =\sum_{k=0}^{p}(-1)^k\frac{p!}{(p-k)!}
 2^{\binom{k+1}{2}}x^{p-k}\mathcal F(x/2^{k+1})+C.        \label{eq:global-poly}
\end{equation}
For $x>0$ and $\alpha\in\C$,
\begin{equation}
 \int_0^xt^\alpha\mathcal F(t)\dd t
 =\sum_{k=0}^{\infty}(-1)^k\fall{\alpha}{k}
 2^{\binom{k+1}{2}}x^{\alpha-k}\mathcal F(x/2^{k+1}),    \label{eq:global-complex}
\end{equation}
with absolute, locally uniform convergence.
\end{theorem}

The signed blocks in \eqref{eq:global-def} therefore possess two complementary
primitive descriptions: a finite scaled-$\mathcal F$ formula and a locally
finite sum of shifted primitives of $\Up$.  Equating them gives blockwise
Thue--Morse identities that may be useful for formalizing the interaction
between the global binary series and integration.

\section{Integrals of the inverse Fabius function}
\label{sec:inverse}

\subsection{The unweighted primitive}
\begin{theorem}[Exact inverse primitive]
\label{thm:inverse-unweighted}
For $0<y<1$,
\begin{equation}
 \boxed{
 \int\InvF(y)\dd y=y\InvF(y)-F(\InvF(y)/2)+C.}            \label{eq:inverse-unweighted}
\end{equation}
The zero-based normalization is obtained with $C=0$.
\end{theorem}

\begin{proof}
Set $a=\InvF(y)$, so $y=F(a)$.  Changing variables $t=F(u)$ gives
\[
 \int_0^y\InvF(t)\dd t=\int_0^a uF'(u)\dd u
 =ay-\int_0^aF(u)\dd u
 =y\InvF(y)-F(\InvF(y)/2).
\]
The last step is \eqref{eq:p0}.
\end{proof}

At $y=1$, \eqref{eq:inverse-unweighted} gives
$\int_0^1\InvF(y)\dd y=1-F(1/2)=1/2$, as also follows from inverse symmetry.

\subsection{Arbitrary complex weights}
For $\Re\beta>-1$ define
\begin{equation}
 K_\beta(y)=\int_0^yt^\beta\InvF(t)\dd t.                 \label{eq:Kbeta-def}
\end{equation}

\begin{theorem}[Weighted inverse reduction]
\label{thm:inverse-weighted}
Let $0<y\le1$ and put $a=\InvF(y)$.  If $\beta\ne-1$ and the integrals are
interpreted on any common local interval avoiding zero, then
\begin{equation}
 \boxed{
 \int y^\beta\InvF(y)\dd y
 =\frac{y^{\beta+1}\InvF(y)}{\beta+1}
 -\frac1{\beta+1}\int^{\InvF(y)}F(x)^{\beta+1}\dd x+C.}  \label{eq:inverse-weighted-local}
\end{equation}
For $\Re\beta>-1$, the zero-based form is
\begin{equation}
 \boxed{
 K_\beta(y)=\frac{y^{\beta+1}\InvF(y)}{\beta+1}
 -\frac1{\beta+1}\int_0^{\InvF(y)}F(x)^{\beta+1}\dd x.} \label{eq:inverse-weighted}
\end{equation}
At the critical exponent $\beta=-1$,
\begin{equation}
 \boxed{
 K_{-1}(y)=\InvF(y)\log y-
 \int_0^{\InvF(y)}\log F(x)\dd x.}                       \label{eq:inverse-critical}
\end{equation}
\end{theorem}

\begin{proof}
With $x=\InvF(y)$,
\[
 K_\beta(y)=\int_0^a xF(x)^\beta F'(x)\dd x.
\]
For $\beta\ne-1$, integrate $x\,\dd(F(x)^{\beta+1})/(\beta+1)$ by parts.  At
$\beta=-1$, integrate $x\,\dd\log F(x)$.  The endpoint boundary
$x\log F(x)\to0$ follows from the repository's small-argument asymptotic (and
already from sufficiently strong flatness bounds).
\end{proof}

For $p\in\N_0$ this becomes
\begin{equation}
 K_p(y)=\frac{y^{p+1}\InvF(y)}{p+1}
 -\frac1{p+1}\int_0^{\InvF(y)}F(x)^{p+1}\dd x.            \label{eq:inverse-integer}
\end{equation}
The reduction is exact, but nonlinear: unlike the forward monomial primitive,
it introduces a power of $F$.  This identifies the precise missing object
needed for a closed inverse calculus.

\subsection{Endpoint integrability threshold}
\begin{proposition}[Sharp threshold]
\label{prop:inverse-threshold}
The normalized integral $K_\beta(y)$ is absolutely convergent at zero exactly
when
\begin{equation}
 \Re\beta\ge-1.                                            \label{eq:inverse-threshold}
\end{equation}
For $\Re\beta<-1$ it diverges, although the local antiderivative
\eqref{eq:inverse-weighted-local} remains valid away from zero.
\end{proposition}

\begin{proof}
Put $t=\e^{-u}$.  Up to a harmless phase,
\[
 \abs{t^\beta\InvF(t)\dd t}
 =\e^{-(\Re\beta+1)u}\InvF(\e^{-u})\dd u.
\]
By \eqref{eq:J-leading}, $\InvF(\e^{-u})$ decays like
$\sqrt u\,\e^{-\sqrt{2(\log2)u}}$, which is subexponential in $u$ but
integrable.  The assertion follows according as the ordinary exponential
factor decays, is absent, or grows.
\end{proof}

\subsection{Complete inverse moments and order statistics}
Let $M_r=\max(X_1,\ldots,X_r)$ for iid copies of $X$.  Since
$\Prob(M_r\le x)=F(x)^r$,
\begin{equation}
 \int_0^1F(x)^r\dd x=\E(1-M_r).                            \label{eq:max-complement}
\end{equation}
Putting $y=1$ in \eqref{eq:inverse-integer} gives
\begin{equation}
 \boxed{
 \int_0^1y^p\InvF(y)\dd y
 =\frac{\E M_{p+1}}{p+1}.}                                \label{eq:inverse-order-stat}
\end{equation}
Thus inverse moments are normalized expected maxima of Fabius random variables.
This probabilistic meaning is complementary to the algebraic inverse-dyadic
and Lambert--Bell analysis in the repository's inverse-frontier report.

At the critical exponent,
\begin{equation}
 \int_0^1\frac{\InvF(y)}y\dd y=-\int_0^1\log F(x)\dd x,   \label{eq:inverse-log-constant}
\end{equation}
an entropy-like constant attached to the distribution function rather than to
its density.

\section{A finite Thue--Morse formula for the nonlinear inverse remainder}
\label{sec:nonlinear-spline}

The integral of $F^r$ in \eqref{eq:inverse-integer} is not reduced by the linear
primitive ladder.  It nevertheless has an exact finite-level representation.
Let
\begin{equation}
 X_N=\sum_{j=1}^{N}2^{-j}U_j
\end{equation}
and let $F_N$ be its distribution function.  Inclusion--exclusion for unequal
box splines gives
\begin{equation}
 \boxed{
 F_N(x)=\frac{2^{N(N+1)/2}}{N!}
 \sum_{m=0}^{2^N-1}(-1)^{\wt(m)}
 \pospart{x-m/2^N}^{N}.}                                  \label{eq:finite-FN}
\end{equation}
Moreover $F_N\to F$ uniformly.

For an integer $r\ge1$, a cutoff $a\in[0,1]$, and a tuple
$\bm m=(m_1,\ldots,m_r)$, put
\begin{equation}
 b(\bm m)=\max_i\frac{m_i}{2^N},\qquad
 \delta_i=b(\bm m)-\frac{m_i}{2^N}.                       \label{eq:b-delta}
\end{equation}
Define
\begin{align}
 Q_{N,r}(a;\bm m)
 &=\mathbf 1_{\{a>b(\bm m)\}}
 \sum_{k_1,\ldots,k_r=0}^{N}
 \frac{\prod_{i=1}^{r}\binom N{k_i}\delta_i^{N-k_i}}
 {1+k_1+\cdots+k_r}                                       \notag\\
 &\qquad\times(a-b(\bm m))^{1+k_1+\cdots+k_r}.           \label{eq:Q-Nr}
\end{align}

\begin{theorem}[Finite Thue--Morse nonlinear integral]
\label{thm:finite-nonlinear}
For every $N,r\in\N$ and $a\in[0,1]$,
\begin{align}
 \int_0^aF_N(x)^r\dd x
 &=\left(\frac{2^{N(N+1)/2}}{N!}\right)^r
 \sum_{m_1,\ldots,m_r=0}^{2^N-1}
 (-1)^{\sum_i\wt(m_i)}Q_{N,r}(a;\bm m).                  \label{eq:finite-nonlinear}\\
 \int_0^aF(x)^r\dd x
 &=\lim_{N\to\infty}\int_0^aF_N(x)^r\dd x.              \label{eq:nonlinear-limit}
\end{align}
Consequently \eqref{eq:inverse-integer} has a completely explicit finite
Thue--Morse spline limit for every $p\in\N_0$.
\end{theorem}

\begin{proof}
Expand the $r$th power of \eqref{eq:finite-FN}.  For a fixed tuple, the product
vanishes below $b=b(\bm m)$.  Above $b$, set $z=x-b$ and write
\[
 \prod_{i=1}^{r}(x-m_i/2^N)^N
 =\prod_{i=1}^{r}(z+\delta_i)^N.
\]
The binomial theorem followed by termwise integration in $z$ gives
\eqref{eq:Q-Nr}.  Uniform convergence $F_N\to F$ and $0\le F_N,F\le1$ give
\eqref{eq:nonlinear-limit} by dominated convergence.
\end{proof}

The formula is exponentially expensive in both $N$ and $r$.  Its value is
structural: it proves that the nonlinear remainder is still governed by the
same finite Thue--Morse geometry, and it offers a starting point for orbit
compression, Walsh transforms, and $q$-binomial acceleration.

\section{Endpoint asymptotics for inverse primitives}
\label{sec:inverse-endpoint}

Let
\begin{equation}
 T=\log(1/y),\qquad L=\log2,\qquad
 a=\sqrt{2L},\qquad C_0=\frac1{\e\sqrt L}.                 \label{eq:endpoint-constants}
\end{equation}
Then \eqref{eq:J-leading} reads
$\InvF(y)\sim C_0\sqrt T\e^{-a\sqrt T}$.

\begin{theorem}[Subcritical weighted inverse primitives]
\label{thm:inverse-endpoint-subcritical}
Let $\rho=\beta+1$ with $\Re\rho>0$.  As $y\downarrow0$,
\begin{equation}
 \boxed{
 K_\beta(y)\sim\frac{y^\rho\InvF(y)}{\rho}.}              \label{eq:Kbeta-leading}
\end{equation}
For real $\rho>0$, the leading model admits the refined endpoint denominator
\begin{equation}
 K_\beta(y)=
 \frac{y^\rho\InvF(y)}{\rho+\dfrac{a}{2\sqrt T}-\dfrac1{2T}}
 \left(1+\bigO(T^{-1})+\text{periodic-transfer terms}\right), \label{eq:Kbeta-refined-structure}
\end{equation}
where a fully rigorous coefficient at the displayed relative order requires
inserting the corpus's periodic next term for $\InvF$.
\end{theorem}

\begin{proof}
After $t=\e^{-u}$,
\[
 K_\beta(y)=\int_T^\infty\e^{-\rho u}\InvF(\e^{-u})\dd u.
\]
The factor $\InvF(\e^{-u})$ is slowly varying on additive $u$-scales compared
with $\e^{-\rho u}$.  Direct one-sided Laplace estimation using
\eqref{eq:J-leading} gives \eqref{eq:Kbeta-leading}.  Differentiating the
leading exponent
$\rho u+a\sqrt u-\frac12\log u$ gives the denominator in
\eqref{eq:Kbeta-refined-structure}; the repository's periodic correction must
be retained before promoting the indicated structure to a complete second-term
theorem.
\end{proof}

At the critical exponent the ordinary exponential disappears and the scale
changes.

\begin{theorem}[Critical inverse primitive]
\label{thm:inverse-endpoint-critical}
As $y\downarrow0$,
\begin{align}
 K_{-1}(y)
 &\sim\frac{\sqrt2}{\e\log2}
 T\exp\!\left(-\sqrt{2(\log2)T}\right)                    \label{eq:Kminus1-leading}\\
 &\sim\sqrt{\frac{2T}{\log2}}\,\InvF(y).                  \label{eq:Kminus1-ratio}
\end{align}
\end{theorem}

\begin{proof}
Now
\[
 K_{-1}(y)=\int_T^\infty\InvF(\e^{-u})\dd u
 \sim C_0\int_T^\infty\sqrt u\,\e^{-a\sqrt u}\dd u.
\]
Set $v=\sqrt u$.  The integral becomes
$2C_0\int_{\sqrt T}^\infty v^2\e^{-av}\dd v$, whose leading
term is $(2C_0/a)T\e^{-a\sqrt T}$.  Substitute the constants.
\end{proof}

For $\Re\beta<-1$, a primitive based at a fixed $y_0>0$ diverges at zero with
leading magnitude $-y^{\beta+1}\InvF(y)/(\beta+1)$.  Thus $\beta=-1$ is a true
transition: above it, the endpoint factor is $1/(\beta+1)$; at it, an extra
$\sqrt{\log(1/y)}$ appears; below it, the same boundary term drives divergence.

\section{Numerical experiments and exact verification}
\label{sec:numerics}

The accompanying script \texttt{experiments.py} is network-free and performs
three logically separate checks.

\begin{enumerate}
\item It computes $V_n=F(2^{-n})$ exactly from \eqref{eq:V-recurrence} using
rational arithmetic.
\item It represents $F(x/2^j)$ as a formal symbol $F_j$ with derivative rule
$F_j'=2^{1-j}F_{j-1}$ and verifies the cancellation in
\eqref{eq:monomial-closed} exactly for $0\le p\le12$.
\item It compares the series \eqref{eq:complex-series} at $x=1$ with a
scrambled Sobol simulation of the random series \eqref{eq:X-random}, using
\eqref{eq:Mellin-prob}.  This last item is only an independent numerical check.
\end{enumerate}

The exact recurrence begins
\begin{equation}
 V_0=1,\quad V_1=\frac12,\quad V_2=\frac5{72},\quad
 V_3=\frac1{288},\quad V_4=\frac{143}{2073600}.            \label{eq:V-first}
\end{equation}
All formal derivative tests and the two moment evaluations
\eqref{eq:moment-dyadic}--\eqref{eq:moment-prob} pass exactly.

\begin{table}[ht]
\centering
\small
\begin{tabular}{@{}rrrr@{}}
\toprule
$\alpha$ & series, 240 terms & Sobol mean & scramble spread \\
\midrule
$-1$   & $0.7582400001750290$ & $0.758239962896570$ & $2.70\times10^{-7}$\\
$-1/2$ & $0.6073863107353257$ & $0.607386310704139$ & $2.51\times10^{-8}$\\
$1/2$  & $0.4209681733402818$ & $0.420968175002922$ & $4.38\times10^{-10}$\\
$2.3$  & $0.2591812728638399$ & $0.259181273425061$ & $5.26\times10^{-11}$\\
\bottomrule
\end{tabular}
\caption{Independent check of $I_\alpha(1)$.  Each Sobol mean uses four
scrambles of $2^{19}$ points and 28 random-series coordinates.  The finite tail
and QMC sampling error explain the displayed differences.}
\label{tab:numerics}
\end{table}

The slower apparent convergence at $\alpha=-1$ is expected: the term quotient
is only polynomial after the superexponential dyadic factors have canceled.
The analytic proof allows an arbitrarily strong polynomial majorant by choosing
a higher flatness order $s$, but the practical constant then worsens.  This is
a natural target for Richardson acceleration based on the exact tail shape.

\section{Further deductions}
\label{sec:further}

\subsection{An exponential-weight generating primitive}
Taking $g(x)=\e^{zx}$ in the infinite version of
\cref{thm:finite-weighted} gives, for $0<x\le1$,
\begin{equation}
 \int_0^x\e^{zt}F(t)\dd t
 =\e^{zx}\sum_{k=0}^{\infty}(-z)^k
 2^{\binom{k+1}{2}}F(x/2^{k+1}).                          \label{eq:exponential-primitive}
\end{equation}
The series is entire in $z$.  Expanding both sides recovers all polynomial
primitives simultaneously.  At $x=1$ it provides a direct dyadic expansion for
the integrated moment-generating function of $F$.

\subsection{A uniqueness principle for the dyadic coefficient row}
Suppose a family of coefficients $c_k$ satisfies
\[
 \int P(x)F(x)\dd x=\sum_{k\ge0}c_kP^{(k)}(x)F(x/2^{k+1})+C
\]
for all polynomials $P$, with only finitely many nonzero terms for each $P$.
Testing successively on $1,x,x^2,\ldots$ forces
$c_k=(-1)^k2^{\binom{k+1}{2}}$.  Thus the triangular coefficient row in
\eqref{eq:polynomial-closed} is the unique local derivative functional built
from the scale tower $F(x/2^{k+1})$.

\subsection{A critical complement contrast}
There is an instructive asymmetry at exponent $-1$:
\begin{equation}
 \int_0^x\frac{F(t)}t\dd t<\infty,
 \qquad
 \int_0^x\frac{1-F(t)}t\dd t=\infty.                      \label{eq:critical-contrast}
\end{equation}
The first has the positive series \eqref{eq:negative-one-series}; the second
has only the local primitive \eqref{eq:complement-minus-one}.  After folding,
the same obstruction is the pole of the fractional up-moment at
$\alpha=-1$.

\section{Conjectures and frontier directions}
\label{sec:conjectures}

\begin{conjecture}[Dyadic-tree algebra for nonlinear primitives]
\label{conj:nonlinear-tree}
For each integer $r\ge2$, the functions
\[
 Q_r(a)=\int_0^aF(x)^r\dd x
\]
admit an absolutely convergent expansion indexed by finite rooted dyadic trees.
The vertex weights should be the same inverse-dyadic moments $d_n$, while tree
symmetries should produce partial Bell polynomials and half-base Gaussian
binomial coefficients.  The finite spline formula
\eqref{eq:finite-nonlinear} is expected to collapse to this tree expansion
after Walsh--Thue--Morse orbit compression.
\end{conjecture}

\begin{conjecture}[Periodic Lambert transfer for inverse primitives]
\label{conj:periodic-transfer}
Let the sharp inverse expansion be written in the repository's phase variable as
\[
 \InvF(y)=\InvF_0(y)\left(1+\frac{A_1(T,\theta)}{\sqrt T}
 +\frac{A_2(T,\theta)}T+\cdots\right),
\]
where the $A_j$ belong to the differential algebra generated by the periodic
Gamma--zeta correction.  Then, for every fixed $\Re\beta>-1$, $K_\beta(y)$ has
an all-orders expansion in the same differential algebra, obtained by the
formal resolvent
\[
 \left(\beta+1+\frac{\dd}{\dd T}\right)^{-1}.
\]
At $\beta=-1$ the resolvent becomes singular and raises the polynomial degree in
$\sqrt T$ by one, explaining \eqref{eq:Kminus1-leading}.
\end{conjecture}

\begin{conjecture}[Sharp denominator law for polynomial primitives]
\label{conj:denominator}
At a dyadic argument $x=a/2^n$, the reduced denominator of the finite primitive
\eqref{eq:monomial-closed} has a predictable $2$-adic valuation obtained by the
minimum of the known valuations of
$F(a/2^{n+k+1})$, except on explicitly classifiable cancellation congruence
classes governed by Thue--Morse parity.  The odd part should factor through the
Mersenne products appearing in the half-base $q$-Pochhammer symbols.
\end{conjecture}

\begin{conjecture}[Stieltjes structure of the negative-power row]
\label{conj:stieltjes}
For each $m\in\N$ and fixed $x\in(0,1]$, the positive terms in
\eqref{eq:negative-series}, after a natural factorial normalization, form a
Stieltjes moment sequence.  If true, diagonal Pad\'e approximants would give
alternating certified bounds and a substantially faster evaluation of
$I_{-m}(x)$.
\end{conjecture}

\begin{conjecture}[Continuous primitive index and dilation cocycle]
\label{conj:fractional-index}
Let $A_\nu=\Vol^\nu F$ be the Riemann--Liouville fractional primitive, which
already supplies the semigroup law $A_{\nu+\mu}=\Vol^\nu A_\mu$.  There is a
canonical transform $\mathscr T_\nu$, holomorphic in $\nu$ for $\Re\nu>0$, and
a scalar cocycle $c(\nu)$ such that
\[
 A_\nu(x)=c(\nu)\,(\mathscr T_\nu F)(2^{-\nu}x),
 \qquad c(n)=2^{\binom n2},\qquad \mathscr T_nF=F
 \quad(n\in\N).
\]
In Mellin coordinates, $c$ should be expressible through a half-base
$q$-gamma or Barnes-type function, while $\mathscr T_\nu$ should become the
identity exactly at integer indices.  The obstruction is that the integer
formula uses the special vector $F$ and does not arise from an operator
identity on arbitrary functions.
\end{conjecture}

\subsection{Promising research program}
The following projects are concrete and separable.

\begin{enumerate}
\item \textbf{Lean formalization.}  Prove \cref{thm:primitive-ladder} first;
then derive the finite polynomial identity by induction.  The complex series
requires an effective flatness lemma, a falling-factorial derivative library,
and a uniform remainder theorem.
\item \textbf{Certified fractional evaluation.}  Combine the exact rational recurrence \eqref{eq:V-recurrence} with an
explicit formal flatness estimate to produce interval bounds for
\eqref{eq:complex-series} at rational $\alpha$.
\item \textbf{$q$-Richardson acceleration.}  Determine the full asymptotic
expansion of the positive tail in \eqref{eq:negative-series}.  The geometric
scale suggests the $q=1/4$ Lagrange filters already developed in the frontier
corpus.
\item \textbf{Nonlinear Walsh compression.}  Replace the $2^{Nr}$ tuple sum in
\eqref{eq:finite-nonlinear} by orbit sums under coordinate permutation,
bitwise XOR, and dyadic reflection; search for Gaussian-binomial collapse.
\item \textbf{Inverse periodic transfer.}  Insert the established periodic
Lambert--Bell expansion of $\InvF$ into the endpoint integrals and prove
\cref{conj:periodic-transfer} with uniform phase derivatives.
\item \textbf{Fourier interpretation.}  Express the even up-moments generated
by \eqref{eq:up-fractional-moment} as fractional derivatives at the origin of
the infinite sinc product, clarifying branch choices and possible Mellin--Barnes
continuation.
\item \textbf{Arithmetic experiments.}  Tabulate reduced denominators of
\eqref{eq:monomial-closed} at general dyadic $x$ and compare them with the
formal $2$-adic lower bounds and Mersenne factors in the current Lean corpus.
\end{enumerate}

\section{Formalization roadmap}
\label{sec:lean-roadmap}

A minimal dependency graph is as follows.

\begingroup
\footnotesize
\begin{longtable}{@{}>{\raggedright\arraybackslash}p{0.32\textwidth}>{\raggedright\arraybackslash}p{0.39\textwidth}>{\raggedright\arraybackslash}p{0.22\textwidth}@{}}
\toprule
Declaration & Mathematical content & Main dependencies \\
\midrule
\endfirsthead
\toprule
Declaration & Mathematical content & Main dependencies \\
\midrule
\endhead
\decl{fabius.iteratedIntegral} &
$\Vol^nF=2^{\binom n2}S^nF$ on $[0,1]$. &
Fabius differential equation; FTC \\
\decl{fabius.integral_pow_mul} &
Finite formula \eqref{eq:monomial-closed}. &
Primitive ladder; repeated integration by parts \\
\decl{fabius.integral_poly_mul} &
Polynomial functional \eqref{eq:polynomial-closed}. &
Linearity; monomial theorem \\
\decl{fabius.complexPower_series} &
Series \eqref{eq:complex-series}. &
Effective flatness; Pochhammer bounds \\
\decl{fabius.mellin_entire} &
Entire Mellin transform and series. &
Complex-power theorem; holomorphic domination \\
\decl{up.integral_pow_mul} &
Piecewise formula \eqref{eq:Hp-piecewise}. &
Reflection; moment identities \\
\decl{invFabius.integral} &
Formula \eqref{eq:inverse-unweighted}. &
Inverse differentiability; primitive $F(x/2)$ \\
\decl{invFabius.weightedIntegral} &
Reduction \eqref{eq:inverse-weighted}. &
Change of variables; power chain rule \\
\decl{invFabius.endpointIntegral} &
Equivalents \eqref{eq:Kbeta-leading}, \eqref{eq:Kminus1-leading}. &
Sharp inverse asymptotic; tail Laplace lemma \\
\bottomrule
\end{longtable}
\endgroup

The names are proposed Lean-style placeholders only.  An implementation should
select repository-consistent declarations after checking the existing
namespace; no claim is made here that any listed declaration already exists.

\section{Conclusion}
\label{sec:conclusion}

The Fabius differential equation contains an integration rule that is simpler
than its differentiation rule: one integration merely halves the argument.
Noncommutation of integration and halving accumulates the triangular factor
$2^{\binom n2}$.  Once this ladder is isolated, every polynomial-weighted
antiderivative is finite, every complex-power antiderivative becomes a
convergent series, and the Mellin transform becomes an entire dyadic
interpolant.  Reflection turns the same calculus into exact formulas for
$1-F(x)$, $F(1-x)$, and $\Up$; the signed extension makes it global.

The inverse function has a different but equally clean structure.  Its
weighted primitives reduce exactly to nonlinear powers of $F$, revealing the
specific obstacle to an elementary inverse calculus.  Those nonlinear powers
still possess a finite Thue--Morse spline limit, and the Lambert--$W$ endpoint
theory transfers to the primitives with a sharp critical transition at weight
$y^{-1}$.

The main practical outcome is a compact toolbox: finite formulas for integer
powers, convergent formulas for arbitrary complex powers, exact arithmetic at
complete or dyadic endpoints, and a clear list of nonlinear and asymptotic
problems remaining for future work.

\appendix
\section{Proof details for compact-uniform convergence}
\label{app:uniform}

Let $K\subset\C$ be compact and let
$\sigma_- =\min_{\alpha\in K}\Re\alpha$.  Choose an integer $s$ with
$s+\sigma_-+1>1$.  By \cref{lem:volterra-flatness}, the absolute value of the
$k$th term of \eqref{eq:complex-series} is bounded by
\begin{equation}
 C_s x^{\sigma_-+s+1}
 \frac{\abs{\fall{\alpha}{k}}}{(s+1)_{k+1}}.              \label{eq:uniform-term-bound}
\end{equation}
On a compact set, the Pochhammer polynomial admits a uniform estimate
\begin{equation}
 \abs{\fall{\alpha}{k}}\le C_K\,\Gamma(k-\sigma_-+c_K)   \label{eq:pochhammer-uniform}
\end{equation}
for a fixed harmless $c_K$.  Gamma-ratio asymptotics then give a summable
majorant $C_{K,s}k^{-s-\sigma_- -2+c_K}$ after increasing $s$ if necessary.
The same calculation with $(s+1)_N$ rather than $(s+1)_{N+1}$ proves uniform
vanishing of the remainder.  Uniformity in $x$ holds on every compact
subinterval $[\varepsilon,1]$; at $x=0$ the primitive is identically zero but
$x^\alpha$ itself has no compact-uniform meaning across arbitrary complex
$\alpha$.

\section{Corpus inventory by document family}
\label{app:corpus}

The following ledger records the families used to establish the nonduplication
boundary.  The public tree is live; names are those visible on 27 August 2026.

\begingroup\small
\begin{longtable}{@{}>{\raggedright\arraybackslash}p{0.30\textwidth}>{\raggedright\arraybackslash}p{0.63\textwidth}@{}}
\toprule
Family & Role in the audit \\
\midrule
\endfirsthead
\toprule
Family & Role in the audit \\
\midrule
\endhead
\familyname{Fabius_Function_and_Rvachev_Up} &
Primary synthesis: definitions, probability, moments, exact inverse-dyadic and
general dyadic values, Bernoulli recurrences, $q$-binomial formulas, global
binary series, Fourier products, asymptotics, and Lean crosswalk. \\
\familyname{non-formalized-research-frontiers} &
Consolidated integration, repeated-integration, dyadic web, inverse,
small-argument, Fourier-decay, and gap-register dossiers. \\
Glossary, non-elementarity, and Lean walkthrough &
Normalization and notation checks; inverse and elementary-function boundaries;
formalization status. \\
Paper transcriptions &
Fabius, Arias de Reyna, arithmetic, lacunary-product, and related source
material used by the repository. \\
Current draft families &
\familyname{Fabius_Arithmetic_Rays_Frontier_Report},
\familyname{Fabius_Dyadic_Self_Sampling_Frontier_Package},
\familyname{Fabius_Half_Integer_Spectral_Frontier_Report},
\familyname{Fabius_Inverse_Frontier_Report_Source_and_PDF},
\familyname{Fabius_Newton_Rvachev_Frontier_Report}, the two
\familyname{Fabius_Rvachev_Frontier_Report} families,
\familyname{Fabius_Rvachev_New_Frontiers},
\familyname{Fabius_Rvachev_Thue_Morse_Frontier_Results},
\familyname{Spectral_Arithmetic_Pascal_Rvachev_Hierarchy},
\familyname{Thue_Morse_Formula_Atlas}, the dyadic-inverse-Barnes,
new-results, results-bundle, spectral-endpoint, and
\familyname{rvachev_up_fourier_decay} families.  These set the newest
nonduplication boundary for inverse, Bell, Mellin, $q$-Richardson, and sinc
claims. \\
Archive manifest &
Frozen early notes; standalone non-elementarity, small-argument, and inverse
dyadic reports; six primary-exposition generations; two supplements; two
$q$-calculus reports; two Lean reports; two glossary versions; and the frozen
dyadic-asymptotic, alternative-representation, formula-to-asymptotic,
$q$-connection, Thue--Morse convergence, gap-register, repeated-integration,
and formula-atlas reports. \\
Auxiliary TeX fragments &
Generated tables and display fragments in self-sampling and numerical bundles;
checked through their parent reports and file inventories rather than treated
as independent mathematical claims. \\
\bottomrule
\end{longtable}
\endgroup

The audit found integration in the senses of moments, probability, quadrature,
and construction by repeated integration.  It did not find the scale-ladder
antiderivative formulas, the exact all-complex-exponent series, or the weighted
inverse reductions of this report.

\section{Reproducibility files}
\label{app:repro}

The accompanying archive contains:
\begin{itemize}
\item \texttt{Fabius\_Monomial\_Antiderivatives.tex}, this source;
\item \texttt{Fabius\_Monomial\_Antiderivatives.pdf}, the rendered report;
\item \texttt{experiments.py}, exact and numerical checks with detailed comments;
\item \texttt{numerical\_results.txt}, the full console report;
\item \texttt{series\_convergence.csv}, selected partial sums;
\item \texttt{README.md} and \texttt{SHA256SUMS}.
\end{itemize}
Run
\begin{lstlisting}[style=pseudo]
python experiments.py --output-dir .
\end{lstlisting}
to regenerate the numerical text and CSV files.  The script requires Python 3,
\texttt{mpmath}, \texttt{numpy}, and \texttt{scipy}; it makes no network
requests.

\begin{thebibliography}{99}

\bibitem{Fabius1966}
J.~Fabius,
\newblock A probabilistic example of a nowhere analytic $C^\infty$-function,
\newblock \emph{Zeitschrift f\"ur Wahrscheinlichkeitstheorie und Verwandte Gebiete}
\textbf{5} (1966), 173--174.
\newblock DOI: \nolinkurl{10.1007/BF00536652}.

\bibitem{Rvachev1990}
V.~A.~Rvachev,
\newblock Compactly supported solutions of functional-differential equations and their applications,
\newblock \emph{Russian Mathematical Surveys} \textbf{45} (1990), no.~1, 87--120.

\bibitem{Arias1982}
J.~Arias de Reyna,
\newblock Definici\'on y estudio de una funci\'on indefinidamente diferenciable de soporte compacto,
\newblock \emph{Revista de la Real Academia de Ciencias Exactas, F\'isicas y Naturales de Madrid}
\textbf{76} (1982), no.~1, 21--38;
English translation, \emph{An infinitely differentiable function with compact support: Definition and properties}, arXiv:\nolinkurl{1702.05442} (2017).

\bibitem{AriasArithmetic}
J.~Arias de Reyna,
\newblock Arithmetic of the Fabius function,
\newblock \emph{INTEGERS} \textbf{18} (2018), Article A51;
arXiv:\nolinkurl{1702.06487}.

\bibitem{Haugland2016}
J.~K.~Haugland,
\newblock Evaluating the Fabius function,
\newblock arXiv:\nolinkurl{1609.07999} (2016).

\bibitem{AlloucheShallit}
J.-P.~Allouche and J.~Shallit,
\newblock \emph{Automatic Sequences: Theory, Applications, Generalizations},
\newblock Cambridge University Press, 2003.

\bibitem{GasperRahman}
G.~Gasper and M.~Rahman,
\newblock \emph{Basic Hypergeometric Series}, second edition,
\newblock Cambridge University Press, 2004.

\bibitem{Comtet}
L.~Comtet,
\newblock \emph{Advanced Combinatorics},
\newblock D.~Reidel, 1974.

\bibitem{Corless1996}
R.~M.~Corless, G.~H.~Gonnet, D.~E.~G.~Hare, D.~J.~Jeffrey, and D.~E.~Knuth,
\newblock On the Lambert $W$ function,
\newblock \emph{Advances in Computational Mathematics} \textbf{5} (1996), 329--359.

\bibitem{ProveItPrimary}
V.~Reshetnikov,
\newblock \emph{The Fabius Function and Rvachev's Up-Function: Exact Arithmetic,
Probability, Fourier Analysis, and Corrected Sharp Asymptotics},
\newblock \texttt{ProveIt}, current source at
\url{https://github.com/VladimirReshetnikov/ProveIt/tree/main/Analysis/FabiusFunction/docs/Fabius_Function_and_Rvachev_Up}.

\bibitem{ProveItFrontiers}
V.~Reshetnikov,
\newblock \emph{Non-formalized Research Frontiers for the Fabius--Rvachev System},
\newblock \texttt{ProveIt}, current source at
\url{https://github.com/VladimirReshetnikov/ProveIt/tree/main/Analysis/FabiusFunction/docs/non-formalized-research-frontiers}.

\bibitem{ProveItInverse}
V.~Reshetnikov,
\newblock \emph{Inverse Frontiers for the Fabius--Rvachev System},
\newblock draft report in \texttt{ProveIt}, 27 August 2026,
\url{https://github.com/VladimirReshetnikov/ProveIt/tree/main/Analysis/FabiusFunction/docs/non-formalized-research-frontiers/drafts/Fabius_Inverse_Frontier_Report_Source_and_PDF}.

\bibitem{ProveItArchive}
V.~Reshetnikov,
\newblock Historical documentation archive and manifest,
\newblock \texttt{ProveIt},
\url{https://github.com/VladimirReshetnikov/ProveIt/tree/main/Analysis/FabiusFunction/docs/archive}.

\end{thebibliography}

\end{document}
